\documentclass[aspectratio=169]{beamer}

\usetheme{Madrid}
\usecolortheme{default}
\usepackage{tikz}
\usepackage{graphicx}
\usepackage{amsmath}

\setbeamertemplate{navigation symbols}{}
\usepackage{xcolor}
\hypersetup{
    colorlinks=true,
    linkcolor=black,
    urlcolor=black,
    citecolor=black
}
\setbeamercolor{bibliography entry author}{fg=black}
\setbeamercolor{bibliography entry title}{fg=black}
\setbeamercolor{bibliography entry location}{fg=black}
\setbeamercolor{bibliography entry note}{fg=black}
\setbeamercolor{bibliography entry}{fg=black}

\definecolor{myskin}{RGB}{245,245,220}   
\setbeamercolor{title}{fg=Black,bg=myskin}
\setbeamercolor{frametitle}{fg=Black,bg=myskin}
\setbeamercolor{structure}{fg=myskin}

% Increase title size
\setbeamerfont{title}{size=\Huge,series=\bfseries}
\setbeamerfont{subtitle}{size=\Large}
\setbeamerfont{author}{size=\Large}
\setbeamerfont{institute}{size=\normalsize}
\setbeamerfont{date}{size=\large}

% Presentation details
\title{Mathematics for Computing}
\subtitle{Assignment -  Module 2}

\author{Anlin Davis}

\institute{M.Tech AI \& SE}
\date{}   % Removes the date
\setbeamertemplate{footline}{}
\begin{document}

\begin{frame}
\titlepage
\end{frame}




%================================================
% SECTION 1
%================================================
\section{Vectors and Coordinate Systems}

%------------------------------------------------
\begin{frame}{Vectors}
\small

A vector is a quantity having both magnitude and direction.

\begin{block}{Representation}
A vector in three dimensions can be written as

$$
\mathbf{v}=
\begin{bmatrix}
v_1\\
v_2\\
v_3
\end{bmatrix}
$$

\end{block}

\begin{itemize}
\item Vectors can represent points, directions and data.
\item In computing, a data sample can be represented as a vector.
\item Example:
\(    \mathbf{x}=
    \begin{bmatrix}
    2\\
    3\\
    5
    \end{bmatrix}
   \)
\end{itemize}

\end{frame}

%------------------------------------------------
\begin{frame}{Coordinate System}
\small

A coordinate system describes the position of a point using numerical coordinates.

\begin{block}{2D Coordinate System}

$$
P=(x,y)
$$

\end{block}

\begin{block}{3D Coordinate System}

$$
P=(x,y,z)
$$

\end{block}

\begin{itemize}
\item The axes are mutually perpendicular.
\item The origin is $(0,0)$ in 2D and $(0,0,0)$ in 3D.
\item Used in computer graphics, robotics and 3D vision.
\end{itemize}

\end{frame}

%------------------------------------------------
\begin{frame}{Vector Addition}
\small

Two vectors can be added component by component.

$$
\mathbf{a}=
\begin{bmatrix}
2\\
3
\end{bmatrix},
\qquad
\mathbf{b}=
\begin{bmatrix}
1\\
4
\end{bmatrix}
$$

$$
\mathbf{a}+\mathbf{b}
=
\begin{bmatrix}
2+1\\
3+4
\end{bmatrix}
=
\begin{bmatrix}
3\\
7
\end{bmatrix}
$$

\begin{itemize}
\item Addition is possible only for vectors of the same dimension.
\item Used for combining directions, positions and feature vectors.
\end{itemize}

\end{frame}

%------------------------------------------------
\begin{frame}{Scalar Multiplication}
\small

A vector can be multiplied by a scalar.

$$
\mathbf{v}=
\begin{bmatrix}
2\\
3
\end{bmatrix},
\qquad
3\mathbf{v}
=
\begin{bmatrix}
6\\
9
\end{bmatrix}
$$

\begin{itemize}
\item Changes the magnitude of a vector.
\item A negative scalar reverses its direction.
\item Used in scaling and normalization.
\end{itemize}

\end{frame}

%================================================
% SECTION 2
%================================================
\section{Linear Combination, Span and Basis}

%------------------------------------------------
\begin{frame}{Linear Combination}
\small

A linear combination is formed by multiplying vectors by scalars and adding them.

$$
\mathbf{v}
=
c_1\mathbf{v}_1+
c_2\mathbf{v}_2+\cdots+c_n\mathbf{v}_n
$$

Example:

$$
\mathbf{v}_1=
\begin{bmatrix}
1\\
0
\end{bmatrix},
\qquad
\mathbf{v}_2=
\begin{bmatrix}
0\\
1
\end{bmatrix}
$$

$$
2\mathbf{v}_1+3\mathbf{v}_2
=
\begin{bmatrix}
2\\
3
\end{bmatrix}
$$

\end{frame}

%------------------------------------------------
\begin{frame}{Span}
\small

The span of a set of vectors is the collection of all their possible
linear combinations.

$$
\operatorname{span}\{\mathbf{v}_1,\mathbf{v}_2\}
=
\{c_1\mathbf{v}_1+c_2\mathbf{v}_2\}
$$

For

$$
\mathbf{e}_1=
\begin{bmatrix}
1\\
0
\end{bmatrix},
\qquad
\mathbf{e}_2=
\begin{bmatrix}
0\\
1
\end{bmatrix}
$$

their span is the entire $\mathbb{R}^2$ plane.

\begin{itemize}
\item Span describes the space generated by vectors.
\item Important in feature representation and dimensionality reduction.
\end{itemize}

\end{frame}

%------------------------------------------------
\begin{frame}{Basis Vectors}
\small

A basis is a set of linearly independent vectors that spans a vector space.

Standard basis of $\mathbb{R}^2$:

$$
\mathbf{e}_1=
\begin{bmatrix}
1\\
0
\end{bmatrix},
\qquad
\mathbf{e}_2=
\begin{bmatrix}
0\\
1
\end{bmatrix}
$$

Every vector in $\mathbb{R}^2$ can be written as

$$
\mathbf{x}=x_1\mathbf{e}_1+x_2\mathbf{e}_2
$$

\begin{block}{Key Idea}
A basis provides a coordinate system for representing vectors.
\end{block}

\end{frame}

%================================================
% SECTION 3
%================================================
\section{Matrices and Transformations}

%------------------------------------------------
\begin{frame}{Matrix Multiplication}
\small

Matrices can be multiplied when their dimensions are compatible.

$$
A=
\begin{bmatrix}
1&2\\
3&4
\end{bmatrix},
\qquad
B=
\begin{bmatrix}
5&6\\
7&8
\end{bmatrix}
$$

$$
AB=
\begin{bmatrix}
19&22\\
43&50
\end{bmatrix}
$$

\begin{itemize}
\item Matrix multiplication is generally not commutative.
\item Matrices represent transformations and data in computing.
\end{itemize}

\end{frame}

%------------------------------------------------
\begin{frame}{Composition of Transformations}
\small

Two transformations can be applied one after another.

$$
T_1(\mathbf{x})=A\mathbf{x},
\qquad
T_2(\mathbf{x})=B\mathbf{x}
$$

Their composition is

$$
T_2(T_1(\mathbf{x}))=BA\mathbf{x}
$$

\begin{block}{Example}
Rotation followed by scaling can be represented by

$$
T(\mathbf{x})=SR\mathbf{x}.
$$

\end{block}

\end{frame}

%------------------------------------------------
\begin{frame}{3D Linear Transformation}
\small

A 3D vector can be transformed using a $3\times3$ matrix.

$$
\mathbf{x}=
\begin{bmatrix}
x\\
y\\
z
\end{bmatrix},
\qquad
\mathbf{x}'=A\mathbf{x}
$$

where

$$
A=
\begin{bmatrix}
a_{11}&a_{12}&a_{13}\\
a_{21}&a_{22}&a_{23}\\
a_{31}&a_{32}&a_{33}
\end{bmatrix}
$$

\begin{itemize}
\item Used in 3D graphics.
\item Used for rotation, scaling and reflection.
\item Important in robotics and computer vision.
\end{itemize}

\end{frame}

%================================================
% SECTION 4
%================================================
\section{Determinant and Inverse}

%------------------------------------------------
\begin{frame}{Determinant}
\small

The determinant is a scalar value associated with a square matrix.

For

$$
A=
\begin{bmatrix}
a&b\\
c&d
\end{bmatrix},
\qquad
\det(A)=ad-bc
$$

Example:

$$
A=
\begin{bmatrix}
2&1\\
3&4
\end{bmatrix}
$$

$$
\det(A)=8-3=5
$$

\begin{itemize}
\item $\det(A)=0$ means the matrix is singular.
\item A non-zero determinant indicates that an inverse exists.
\end{itemize}

\end{frame}

%------------------------------------------------
\begin{frame}{Inverse of a Matrix}
\small

For a $2\times2$ matrix,

$$
A=
\begin{bmatrix}
a&b\\
c&d
\end{bmatrix}
$$

the inverse is

$$
A^{-1}
=
\frac{1}{ad-bc}
\begin{bmatrix}
d&-b\\
-c&a
\end{bmatrix}
$$

provided that

$$
ad-bc\neq0.
$$

\begin{block}{Important}

$$
AA^{-1}=A^{-1}A=I
$$

\end{block}

\end{frame}

%------------------------------------------------
\begin{frame}{Methods for Finding the Inverse}
\small

Common methods include:

\begin{enumerate}
\item Formula method for small matrices.
\item Gaussian elimination.
\item Gauss--Jordan elimination.
\item Adjoint and determinant method.
\end{enumerate}

For Gauss--Jordan elimination,

$$
[A\mid I]\longrightarrow[I\mid A^{-1}]
$$

\begin{itemize}
\item Useful for solving systems of linear equations.
\item Widely used in numerical computing.
\end{itemize}

\end{frame}

%================================================
% SECTION 5
%================================================
\section{Column Space and Null Space}

%------------------------------------------------
\begin{frame}{Column Space}
\small

The column space of a matrix is the span of its column vectors.

For

$$
A=
\begin{bmatrix}
1&2\\
2&4
\end{bmatrix}
$$

the columns are

$$
\begin{bmatrix}
1\\
2
\end{bmatrix},
\qquad
\begin{bmatrix}
2\\
4
\end{bmatrix}
$$

Therefore,

$$
\operatorname{Col}(A)
=
\operatorname{span}
\left\{
\begin{bmatrix}
1\\
2
\end{bmatrix}
\right\}
$$

\begin{itemize}
\item The dimension of the column space is the rank.
\end{itemize}

\end{frame}

%------------------------------------------------
\begin{frame}{Null Space}
\small

The null space contains all vectors satisfying

$$
A\mathbf{x}=\mathbf{0}.
$$

For

$$
A=
\begin{bmatrix}
1&2\\
2&4
\end{bmatrix}
$$

we solve

$$
x+2y=0
\qquad\Rightarrow\qquad
x=-2y
$$

Thus,

$$
\mathbf{x}
=
y
\begin{bmatrix}
-2\\
1
\end{bmatrix}
$$

Therefore,

$$
\operatorname{Null}(A)
=
\operatorname{span}
\left\{
\begin{bmatrix}
-2\\
1
\end{bmatrix}
\right\}
$$

\end{frame}

%================================================
% SECTION 6
%================================================
\section{Norms and Square Matrices}

%------------------------------------------------
\begin{frame}{Vector Norm}
\small

A norm measures the size or length of a vector.

For

$$
\mathbf{x}=
\begin{bmatrix}
x_1\\
x_2\\
\vdots\\
x_n
\end{bmatrix}
$$

the Euclidean norm is

$$
\|\mathbf{x}\|_2
=
\sqrt{x_1^2+x_2^2+\cdots+x_n^2}
$$

Example:

$$
\left\|
\begin{bmatrix}
3\\
4
\end{bmatrix}
\right\|_2
=5
$$

\begin{itemize}
\item Used in distance calculation and Important in machine learning and optimization .

\end{itemize}

\end{frame}

%------------------------------------------------
\begin{frame}{Square Matrices}
\small

A square matrix has the same number of rows and columns.

$$
A\in\mathbb{R}^{n\times n}
$$

Example:

$$
A=
\begin{bmatrix}
1&2\\
3&4
\end{bmatrix}
$$

Important properties include:

\begin{itemize}
\item Determinant
\item Inverse
\item Eigenvalues
\item Trace
\item Rank
\end{itemize}

Square matrices are widely used for linear transformations.

\end{frame}

%================================================
% SECTION 7
%================================================
\section{Transformations Between Dimensions}

%------------------------------------------------
\begin{frame}{Transformation Between Dimensions}
\small

A linear transformation maps vectors from one vector space to another.

$$
T:\mathbb{R}^n\rightarrow\mathbb{R}^m
$$

Example:

$$
T:\mathbb{R}^2\rightarrow\mathbb{R}^3
$$

can be represented using a $3\times2$ matrix:

$$
A=
\begin{bmatrix}
1&0\\
0&1\\
1&1
\end{bmatrix}
$$

Then

$$
\mathbf{y}=A\mathbf{x}
$$

\end{frame}

%------------------------------------------------
\begin{frame}{Linear Transformation}
\small

A transformation $T$ is linear if

$$
T(\mathbf{u}+\mathbf{v})
=
T(\mathbf{u})+T(\mathbf{v})
$$

and

$$
T(c\mathbf{u})=cT(\mathbf{u})
$$

\begin{block}{Example}

$$
T(x,y)=(2x,3y)
$$

is a linear transformation.
\end{block}

Linear transformations preserve vector addition and scalar multiplication.

\end{frame}

%================================================
% SECTION 8
%================================================
\section{Dot Product and Duality}

%------------------------------------------------
\begin{frame}{Dot Product}
\small

The dot product of two vectors is

$$
\mathbf{u}\cdot\mathbf{v}
=
u_1v_1+u_2v_2+\cdots+u_nv_n
$$

Example:

$$
\begin{bmatrix}
1\\
2\\
3
\end{bmatrix}
\cdot
\begin{bmatrix}
4\\
5\\
6
\end{bmatrix}
=
4+10+18=32
$$

Geometrically,

$$
\mathbf{u}\cdot\mathbf{v}
=
\|\mathbf{u}\|\|\mathbf{v}\|\cos\theta
$$

\end{frame}

%------------------------------------------------
\begin{frame}{Dot Product and Orthogonality}
\small

Two non-zero vectors are orthogonal when

$$
\mathbf{u}\cdot\mathbf{v}=0
$$

Example:

$$
\begin{bmatrix}
1\\
0
\end{bmatrix}
\cdot
\begin{bmatrix}
0\\
1
\end{bmatrix}
=0
$$

Therefore, the vectors are perpendicular.

\begin{itemize}
\item Used in projections.
\item Used in similarity calculations.
\item Important in machine learning and computer graphics.
\end{itemize}

\end{frame}

%------------------------------------------------
\begin{frame}{Duality}
\small

A vector can define a linear functional through the dot product.

For a fixed vector $\mathbf{a}$,

$$
f(\mathbf{x})=\mathbf{a}^{T}\mathbf{x}
$$

Example:

$$
\mathbf{a}=
\begin{bmatrix}
2\\
3
\end{bmatrix}
$$

gives

$$
f(\mathbf{x})=2x_1+3x_2
$$

Thus, vectors can be used to define linear functions.

\end{frame}

%================================================
% SECTION 9
%================================================
\section{Cross Product}

%------------------------------------------------
\begin{frame}{Cross Product}
\small

The cross product is defined for vectors in $\mathbb{R}^3$.

$$
\mathbf{u}\times\mathbf{v}
=
\begin{bmatrix}
u_2v_3-u_3v_2\\
u_3v_1-u_1v_3\\
u_1v_2-u_2v_1
\end{bmatrix}
$$

Example:

$$
\begin{bmatrix}
1\\
0\\
0
\end{bmatrix}
\times
\begin{bmatrix}
0\\
1\\
0
\end{bmatrix}
=
\begin{bmatrix}
0\\
0\\
1
\end{bmatrix}
$$

\begin{itemize}
\item The result is perpendicular to both vectors.
\end{itemize}

\end{frame}

%------------------------------------------------
\begin{frame}{Cross Product as a Linear Transformation}
\footnotesize

For a fixed vector $\mathbf{a}$,

$$
T(\mathbf{x})=\mathbf{a}\times\mathbf{x}
$$

is a linear transformation.

It can be represented by a skew-symmetric matrix:

$$
[\mathbf{a}]_\times=
\begin{bmatrix}
0&-a_3&a_2\\
a_3&0&-a_1\\
-a_2&a_1&0
\end{bmatrix}
$$

Therefore,

$$
\mathbf{a}\times\mathbf{x}
=
[\mathbf{a}]_\times\mathbf{x}
$$

This representation is useful in robotics and 3D computer vision.

\end{frame}

%================================================
% SECTION 10
%================================================
\section{Kronecker Product}

%------------------------------------------------
\begin{frame}{Kronecker Rule / Product}
\small

The Kronecker product combines two matrices to form a larger matrix.

For

$$
A=
\begin{bmatrix}
a&b\\
c&d
\end{bmatrix}
$$

and matrix $B$,

$$
A\otimes B
=
\begin{bmatrix}
aB&bB\\
cB&dB
\end{bmatrix}
$$

\begin{itemize}
\item The result contains blocks of scaled copies of $B$.
\item Useful for large structured matrix operations.
\item Applications include signal processing and machine learning.
\end{itemize}

\end{frame}

%================================================
% SECTION 11
%================================================
\section{Change of Basis}

%------------------------------------------------
\begin{frame}{Change of Basis}
\small

The same vector can have different coordinate representations under different bases.

Suppose

$$
B=\{\mathbf{b}_1,\mathbf{b}_2\}
$$

is a new basis.

The vector can be written as

$$
\mathbf{x}
=
c_1\mathbf{b}_1+c_2\mathbf{b}_2
$$

The coordinate vector is

$$
[\mathbf{x}]_B=
\begin{bmatrix}
c_1\\
c_2
\end{bmatrix}
$$

\begin{block}{Key Idea}
The vector remains the same; only its coordinates change.
\end{block}

\end{frame}

%================================================
% SECTION 12
%================================================
\section{Eigenvalues and Eigenvectors}

%------------------------------------------------
\begin{frame}{Eigenvectors}
\small

An eigenvector of a matrix $A$ is a non-zero vector whose direction remains unchanged after transformation.

$$
A\mathbf{v}=\lambda\mathbf{v}
$$

where

\begin{itemize}
\item $\mathbf{v}$ is the eigenvector.
\item $\lambda$ is the eigenvalue.
\end{itemize}

Example:

$$
A=
\begin{bmatrix}
2&0\\
0&3
\end{bmatrix}
$$

has eigenvectors

$$
\begin{bmatrix}
1\\
0
\end{bmatrix},
\qquad
\begin{bmatrix}
0\\
1
\end{bmatrix}
$$

\end{frame}

%------------------------------------------------
\begin{frame}{Finding Eigenvalues}
\small

Eigenvalues are found using

$$
\det(A-\lambda I)=0
$$

For

$$
A=
\begin{bmatrix}
2&0\\
0&3
\end{bmatrix}
$$

$$
\det
\begin{bmatrix}
2-\lambda&0\\
0&3-\lambda
\end{bmatrix}
=0
$$

Therefore,

$$
(2-\lambda)(3-\lambda)=0
$$

and

$$
\boxed{\lambda_1=2,\qquad\lambda_2=3}
$$

\end{frame}

%------------------------------------------------
\begin{frame}{Applications of Eigenvalues}
\small

Eigenvalues and eigenvectors are important in computing.

\begin{itemize}
\item Principal Component Analysis (PCA)
\item Dimensionality reduction
\item Image processing
\item Graph analysis
\item Stability analysis
\item Spectral clustering
\end{itemize}

\begin{block}{AI Connection}
PCA uses eigenvectors of a covariance matrix to find important directions in data.
\end{block}

\end{frame}

%================================================
% SECTION 13
%================================================
\section{Vector Spaces}

%------------------------------------------------
\begin{frame}{Vector Spaces}
\small

A vector space is a set of vectors closed under vector addition and scalar multiplication.

For vectors $\mathbf{u}$ and $\mathbf{v}$ in a vector space,

$$
\mathbf{u}+\mathbf{v}\in V
$$

and for scalar $c$,

$$
c\mathbf{u}\in V
$$

Examples include:

\begin{itemize}
\item $\mathbb{R}^2$
\item $\mathbb{R}^3$
\item Polynomial spaces
\item Spaces of matrices
\item Feature spaces in machine learning
\end{itemize}

\end{frame}

%------------------------------------------------
\begin{frame}{Vector Space: Example}
\small

Consider

$$
V=\mathbb{R}^2
$$

Let

$$
\mathbf{u}=
\begin{bmatrix}
1\\
2
\end{bmatrix},
\qquad
\mathbf{v}=
\begin{bmatrix}
3\\
4
\end{bmatrix}
$$

Then

$$
\mathbf{u}+\mathbf{v}
=
\begin{bmatrix}
4\\
6
\end{bmatrix}
\in V
$$

For scalar $2$,

$$
2\mathbf{u}
=
\begin{bmatrix}
2\\
4
\end{bmatrix}
\in V
$$

Thus, $\mathbb{R}^2$ is a vector space.

\end{frame}

%================================================
% SECTION 14
%================================================
\section{Applications in Computing}

%------------------------------------------------
\begin{frame}{Applications in AI and Computing}
\small

Linear algebra is fundamental to modern computing.

\begin{itemize}
\item \textbf{Machine Learning:}
Data and model parameters are represented as vectors and matrices.

```
\item \textbf{Computer Vision:}
Images are represented as matrices.

\item \textbf{Deep Learning:}
Neural networks perform matrix and vector operations.

\item \textbf{Computer Graphics:}
Transformations represent rotation, scaling and projection.

\item \textbf{Robotics:}
Vectors and matrices describe position and orientation.
```

\end{itemize}

\end{frame}


%================================================
% THANK YOU
%================================================

\begin{frame}[plain]

\vfill

\begin{center}
{\LARGE\bfseries Thank You}
\end{center}

\vfill

\end{frame}

%================================================
\end{document}
%================================================
